\documentclass[12pt]{article}
\usepackage{takeaway}
\setband{Grades 4--5}

\newcommand{\fillbox}[1][1.2cm]{\tikz[baseline=-0.6ex]\node[draw=ink,line width=0.7pt,rounded corners=1mm,minimum width=#1,minimum height=0.72cm]{};}
\newcommand{\rul}[1]{{\subheadfont\{#1\}}}

% labelled strip of boxes; #1 from, #2 to, #3 box width (cm); optional prefilled first label #4
\newcommand{\labelstrip}[4]{%
  \begin{tikzpicture}[x=1cm,y=1cm]
  \foreach \n [count=\c from 0] in {#1,...,#2}{%
    \draw[grid,line width=0.6pt,fill=white] (\c*#3,0) rectangle ++(#3,-0.85);
    \node[font=\subheadfont\footnotesize,text=teal] at (\c*#3+0.5*#3,0.25) {\n};}
  \node[font=\headfont\large,text=plum] at (0.5*#3,-0.425) {#4};
  \end{tikzpicture}}

\begin{document}
\fontsize{12}{16}\selectfont

%=====================================================================
\bigtitle{Take-Away Games}{Patterns that repeat, and why they must}
\nameline

\vspace{4pt}
\begin{rulebox}
\begin{itemize}[leftmargin=1.2em,itemsep=0pt,topsep=1pt]
\item Two players share one pile of counters and take turns. On your turn you must take some counters;
the {\subheadfont rule} says how many you may take. You can't take more than the pile has.
\item Whoever takes the {\subheadfont last counter wins}.
\item We write a rule in braces: \rul{1, 2} means ``take 1 or take 2.''
\end{itemize}
{\small\color{grid} Group of three? Two play and the third is the referee who checks every move. Rotate each game.}
\end{rulebox}

\begin{ideabox}[W piles and L piles]
Always look at a pile from the point of view of the player {\subheadfont about to move}.
\begin{itemize}[leftmargin=1.2em,itemsep=0pt,topsep=1pt]
\item {\subheadfont W pile:} the player to move can force a win, whatever the opponent does.
\item {\subheadfont L pile} (a trap): the player to move will lose if the opponent plays well.
\end{itemize}
{\subheadfont Working backward.} 0 is L (your opponent just took the last counter).
A pile is {\subheadfont W} if {\subheadfont at least one} move leads to an L pile.
A pile is {\subheadfont L} if {\subheadfont every} move leads to a W pile.
\end{ideabox}

{\subheadfont\large\color{orange} Warm-up rule: \{1, 2\}}

\prob Play a few games starting with 10 or 11 counters. Who wins, the first or the second player?

\prob Label each pile W or L for the rule \rul{1, 2}.

\begin{center}\labelstrip{0}{15}{1.1}{L}\end{center}

\prob Every multiple of 3 is an L pile. Explain why. Your explanation must work against {\subheadfont every}
move your opponent could make. (What do you do after they take 1? After they take 2?)

\writelines{3}

\prob You face a pile that is {\subheadfont not} a multiple of 3. What is your winning move?

\writelines{1}

\newpage
%=====================================================================
\pagetitle{A new rule: \{1, 3, 4\}}
Take 1, 3, or 4 counters (not 2). With 2 counters left, the only move is to take 1.

\prob Play a few games starting with 10 to 15 counters. Make a guess: which piles do you {\subheadfont not} want on your turn?

\writelines{1}

\prob Label every pile from 0 to 29 as W or L. Work from small to large, using the rules for working backward.

\begin{center}
\begin{tikzpicture}[x=1cm,y=1cm]
  \foreach \r in {0,1,2}{\foreach \c in {0,...,9}{%
    \pgfmathtruncatemacro{\n}{10*\r+\c}
    \draw[grid,line width=0.6pt,fill=white] (\c*1.7,-\r*1.75) rectangle ++(1.6,-1.0);
    \node[font=\subheadfont\footnotesize,text=teal] at (\c*1.7+0.8,-\r*1.75+0.27) {\n};}}
  \node[font=\headfont\large,text=plum] at (0.8,-0.5) {L};
\end{tikzpicture}
\end{center}

\prob List the L piles you found: \blank[11.2cm]

\prob Describe the pattern. Can you say it using remainders?\\[2pt]
A pile is L exactly when dividing it by \fillbox[0.9cm] leaves a remainder of \fillbox[0.9cm] or \fillbox[0.9cm].

\prob Predict W or L, and explain how you know.\\[4pt]
50: \fillbox\quad 100: \fillbox\quad 1000: \fillbox\quad \blank[7.9cm]

\prob The pile has 1000 counters and it is your turn. How many should you take? Why?

\writelines{2}

\newpage
%=====================================================================
\pagetitle{Why does the pattern repeat?}
To label pile \textit{n} under \rul{1, 3, 4}, you only need the labels of piles \textit{n}\,−\,1, \textit{n}\,−\,3, and \textit{n}\,−\,4.
All three sit inside the {\subheadfont window} of 4 piles just before \textit{n}.

\begin{center}
\begin{tikzpicture}[x=1cm,y=1cm]
  \foreach \t [count=\i from 0] in {\textit{n}\,−\,6,\textit{n}\,−\,5,\textit{n}\,−\,4,\textit{n}\,−\,3,\textit{n}\,−\,2,\textit{n}\,−\,1,\textit{n}}{%
    \draw[grid,line width=0.7pt,fill=white] (\i*1.9,0) rectangle ++(1.6,-1.0);
    \node[font=\subheadfont\small,text=teal] at (\i*1.9+0.8,0.3) {\t};}
  \fill[tealsoft] (2*1.9-0.12,-1.15) rectangle (5*1.9+1.72,0.62);
  \foreach \t [count=\i from 0] in {\textit{n}\,−\,6,\textit{n}\,−\,5,\textit{n}\,−\,4,\textit{n}\,−\,3,\textit{n}\,−\,2,\textit{n}\,−\,1,\textit{n}}{%
    \draw[grid,line width=0.7pt,fill=white] (\i*1.9,0) rectangle ++(1.6,-1.0);
    \node[font=\subheadfont\small,text=teal] at (\i*1.9+0.8,0.3) {\t};}
  \draw[teal,line width=1.4pt,rounded corners=2mm] (2*1.9-0.12,-1.15) rectangle (5*1.9+1.72,0.62);
  \node[font=\subheadfont\small,text=teal] at (3.5*1.9+0.8,0.95) {the window};
  \draw[ink,line width=1.6pt] (6*1.9,0) rectangle ++(1.6,-1.0);
  \node[font=\headfont\Large] at (6*1.9+0.8,-0.5) {?};
  \foreach \k/\d/\o in {1/0.55/-0.4,3/1.3/0,4/2.1/0.4}{%
    \pgfmathsetmacro{\xa}{6*1.9+0.8+\o}\pgfmathsetmacro{\xb}{6*1.9-\k*1.9+0.95}
    \draw[-{Stealth[length=3mm]},orange,line width=1.2pt] (\xa,-1.05) .. controls (\xa,-1.05-\d) and (\xb,-1.05-\d) .. (\xb,-1.05);
    \node[anchor=east,inner sep=1pt,font=\footnotesize\subheadfont,text=orange] at (\xb-0.1,-1.4) {take \k};}
  \pgfresetboundingbox
  \path (-0.1,1.2) rectangle (13.1,-2.75);
\end{tikzpicture}
\end{center}

\prob Copy your labels from page 2 into these two windows.

\begin{center}
\begin{tikzpicture}[x=1cm,y=1cm]
  \foreach \s/\x in {0/0,7/9.4}{%
    \foreach \i in {0,...,3}{%
      \pgfmathtruncatemacro{\n}{\s+\i}
      \draw[grid,line width=0.7pt,fill=white] (\x+\i*1.5,0) rectangle ++(1.3,-0.9);
      \node[font=\subheadfont\footnotesize,text=teal] at (\x+\i*1.5+0.65,0.25) {\n};}
    \draw[teal,line width=1.2pt,rounded corners=2mm] (\x-0.15,-1.05) rectangle (\x+4*1.5-0.05,0.55);}
  \node[font=\subheadfont\small] at (2.85,-1.45) {piles 0 to 3};
  \node[font=\subheadfont\small] at (12.25,-1.45) {piles 7 to 10};
\end{tikzpicture}
\end{center}
What do you notice? \blank[13.2cm]

\prob Suppose two windows have exactly the same labels. Explain why the label just after each window must be
the same too, and then why the labels keep matching forever.

\writelines{2}

\prob A window has 4 labels, each W or L. How many different windows are possible?\quad\fillbox

\begin{ideabox}[The pigeonhole principle]
If more pigeons than holes fly into the holes, then some hole gets at least two pigeons.
\end{ideabox}

\prob Explain why {\subheadfont every} rule whose biggest move is 4 must have a win/lose pattern that
eventually repeats. Hint: think about the windows that start at piles 0, 1, 2, \dots, 16.

\writelines{3}

\prob If the biggest move is 10, how many different windows are possible?\quad\fillbox[1.9cm]

\newpage
%=====================================================================
\pagetitle{Change the rule}
\prob Label piles 0 to 19 as W or L for each rule. Then write how long the repeating block is.

\begin{center}
\begin{tikzpicture}[x=1cm,y=1cm]
  \foreach \r/\name in {0/{\{1, 2, 3\}},1/{\{1, 4\}},2/{\{1, 3, 5\}},3/{\{1, 2, 4\}},4/{\{\ \ \ \ \ \ \ \ \ \ \}}}{%
    \begin{scope}[yshift=-\r*2.05cm]
    \node[font=\subheadfont,anchor=west] at (-0.1,-0.42) {\name};
    \foreach \n in {0,...,19}{%
      \draw[grid,line width=0.6pt,fill=white] (2.2+\n*0.79,0) rectangle ++(0.72,-0.85);
      \node[font=\subheadfont\scriptsize,text=teal] at (2.2+\n*0.79+0.36,0.22) {\n};}
    \node[font=\headfont,text=plum] at (2.56,-0.425) {L};
    \node[font=\small,anchor=east] at (16.4,-1.25) {repeating block length:};
    \draw[ink,line width=0.7pt,rounded corners=1mm] (16.55,-0.98) rectangle ++(1.0,-0.6);
    \end{scope}}
  \node[font=\footnotesize\itshape,text=grid,anchor=west] at (-0.1,-8.2+0.3) {your rule:};
\end{tikzpicture}
\end{center}

\prob Under \rul{1, 3, 5} the L piles are exactly the even numbers. Explain why.
(Hint: every move takes an odd number.)

\writelines{2}

\prob \rul{1, 2} and \rul{1, 2, 4} have the same L piles. Why doesn't the extra move ``take 4'' help anyone?

\writelines{2}

\prob {\subheadfont Greedy Gus} always takes the biggest number the rule allows. Under \rul{1, 3, 4}, find a pile
where Gus's move loses, but a different move would have won.\\[4pt]
Pile: \fillbox\quad Gus takes \fillbox\quad A winning move is to take \fillbox

\newpage
%=====================================================================
\pagetitle{Challenge: two piles}
\begin{rulebox}[Two-pile rule]
There are {\subheadfont two} piles. On your turn, choose {\subheadfont one} pile and take 1 or 2 counters from it.
Whoever takes the very last counter wins.
\end{rulebox}

\prob Label each position W or L. The cell in row \textit{a} and column \textit{b} means ``\textit{a} counters in one pile and
\textit{b} in the other.'' A move goes 1 or 2 cells up, or 1 or 2 cells left (the arrows show the moves from the shaded cell, which means piles of 4 and 5).

\begin{center}
\begin{tikzpicture}[x=1cm,y=1cm]
  \def\s{1.25}
  \foreach \a in {0,...,6}{%
    \node[font=\subheadfont,text=teal] at (-0.55,-\a*\s-0.5*\s) {\a};
    \node[font=\subheadfont,text=teal] at (\a*\s+0.5*\s,0.45) {\a};
    \foreach \b in {0,...,6}{%
      \draw[grid,line width=0.6pt,fill=white] (\b*\s,-\a*\s) rectangle ++(\s,-\s);}}
  \node[font=\subheadfont\small,text=teal,anchor=south] at (3.5*\s,0.75) {pile B};
  \node[font=\subheadfont\small,text=teal,rotate=90,anchor=south] at (-0.95,-3.5*\s) {pile A};
  \node[font=\headfont\large,text=plum] at (0.5*\s,-0.5*\s) {L};
  % example moves from (a,b)=(4,5)
  \fill[plumsoft] (5*\s+0.04,-4*\s-0.04) rectangle ++(\s-0.08,-\s+0.08);
  \fill[plum] (5.5*\s,-4.5*\s) circle (0.09);
  \foreach \tx/\ty in {4.5/4.5,3.5/4.5,5.5/3.5,5.5/2.5}{\draw[plum,line width=1pt,fill=white] (\tx*\s,-\ty*\s) circle (0.09);}
  \draw[-{Stealth[length=2.3mm]},plum,line width=1.1pt,shorten <=2pt,shorten >=3pt] (5.5*\s,-4.5*\s) -- (4.5*\s,-4.5*\s);
  \draw[-{Stealth[length=2.3mm]},plum,line width=1.1pt,shorten <=2pt,shorten >=3pt] (5.5*\s,-4.5*\s) to[bend left=35] (3.5*\s,-4.5*\s);
  \draw[-{Stealth[length=2.3mm]},plum,line width=1.1pt,shorten <=2pt,shorten >=3pt] (5.5*\s,-4.5*\s) -- (5.5*\s,-3.5*\s);
  \draw[-{Stealth[length=2.3mm]},plum,line width=1.1pt,shorten <=2pt,shorten >=3pt] (5.5*\s,-4.5*\s) to[bend right=35] (5.5*\s,-2.5*\s);
\end{tikzpicture}
\end{center}

\prob Which positions are L? Describe the pattern. \blank[8.9cm]

\writelines{1}

\prob Explain why two {\subheadfont equal} piles (like 5 and 5, or 100 and 100) are always L. (Hint: copycat.)

\writelines{2}

\prob Without drawing a bigger grid: is the position with piles 14 and 20 W or L? \fillbox[1.4cm]

\newpage
%=====================================================================
\pagetitle{Challenge: beyond win and lose}
Mathematicians give each pile a number called its {\subheadfont Grundy value}. It tells you more than just W or L.

\begin{ideabox}[Grundy values]
The {\subheadfont mex} of a list of numbers is the {\subheadfont smallest} of 0, 1, 2, 3, \dots{} that is
{\subheadfont not} in the list. For example, mex of (0, 1, 3) is 2, and mex of (1, 2) is 0.\\[3pt]
The {\subheadfont Grundy value} of a pile is the mex of the values of all the piles you can move to.
A pile of 0 has value 0. A pile is L exactly when its value is 0.
\end{ideabox}

\prob Find the Grundy values for the rule \rul{1, 3, 4}. Two are done for you. What pattern do you see?

\begin{center}
\begin{tikzpicture}[x=1cm,y=1cm]
  \foreach \n in {0,...,14}{%
    \draw[grid,line width=0.6pt,fill=white] (\n*1.2,0) rectangle ++(1.1,-0.95);
    \node[font=\subheadfont\footnotesize,text=teal] at (\n*1.2+0.55,0.25) {\n};}
  \node[font=\headfont\large,text=plum] at (0.55,-0.475) {0};
  \node[font=\headfont\large,text=plum] at (1.75,-0.475) {1};
\end{tikzpicture}
\end{center}
{\small\color{grid} Example: from pile 1 you can only move to pile 0 (value 0). The mex of (0) is 1.}

Pattern: \blank[15.6cm]

\begin{ideabox}[A fact about two piles]
Play two piles that both use the same rule; on your turn, take from one pile only.
The player about to move is in an {\subheadfont L} position exactly when the two piles have {\subheadfont equal} Grundy values.
\end{ideabox}

\prob Under \rul{1, 3, 4}, piles of 4 and of 6 are both W piles on their own. What are their Grundy values?
What happens if you play with {\subheadfont two} piles, 4 and 6? Check by playing!

\writelines{2}

\prob Find a pile \textit{n} that is not 5 such that two piles, 5 and \textit{n}, make an L position.\quad\fillbox

\begin{challengebox}[Still unsolved]
In some take-away games you may also split a pile into two piles. For a big family of games like these,
called finite octal games, the mathematician Richard Guy guessed that the pattern of Grundy values always repeats in the end.
Nobody has proved it yet, and nobody has found a game where it fails. It is a real open question!\\[3pt]
{\small\color{grid} To read more: Berlekamp, Conway, and Guy, {\itshape Winning Ways}; Siegel, {\itshape Combinatorial Game Theory}.}
\end{challengebox}

\end{document}
